% Part: computability
% Chapter: computability-theory
% Section: s-m-n

\documentclass[../../../include/open-logic-section]{subfiles}

\begin{document}

% SEGMENT OLP-0232-S01

\olfileid{cmp}{thy}{smn}
\olsection{$s$-$m$-$n$ தேற்றம்}

\begin{explain}
அடுத்த தேற்றம் இன்னும் ஒரு கணத்தில் தெளிவாகும் காரணத்தால்
“$s$-$m$-$n$ தேற்றம்” என அறியப்படுகிறது. தேற்றம் என்ன கூறுகிறது
என்பதைப் புரிந்துகொள்வதே கடினமான பகுதி; கூற்றைப் புரிந்துகொண்டதும்
அது ஓரளவு தெளிவாகத் தோன்றும்.
\end{explain}

\begin{thm}
\ollabel{thm:s-m-n}
  இயல் எண்கள் $n$ மற்றும்~$m$ இன் ஒவ்வொரு சோடிக்கும், ஒவ்வொரு
  தொடர் $e$, $a_0$, \dots, $a_{m-1}$, $y_0$ ,\dots,
  $y_{n-1}$ க்கும் பின்வருவது பொருந்துமாறு ஒரு முதன்மை
  மீள்வரையறைச் சார்பு~$s^m_n$ உள்ளது:
  \[
  \cfind{s^m_n(e, a_0, \dots, a_{m-1})}[n](y_0, \dots, y_{n-1}) \simeq
  \cfind{e}[m+n](a_0, \dots, a_{m-1}, y_0, \dots, y_{n-1}).
  \]
\end{thm}

\begin{explain}
$s^m_n$ ஆனது \emph{நிரல்கள்மீது} செயல்படுவதாகக் கருதுவது உதவும்.
அதாவது $(m+n)$-இடச் சார்புக்கான நிரல்~$e$ ஒன்றையும், நிலையான
உள்ளீடுகள் $a_0$, \dots, $a_{m-1}$ ஆகியவற்றையும் $s^m_n$
பெற்றுக்கொள்கிறது; மீதமுள்ள உள்ளீடுகளின் $n$-இடச் சார்புக்கான
நிரல் $s^m_n(x, a_0, \dots, a_{m-1})$ ஐத் திருப்பித்தருகிறது.
\iftag{TMs}{$x$ ஐ ஒரு டியூரிங் பொறியின் விளக்கம் என்று கருதினால்,
$s^m_n(e, a_0, \dots, a_{m-1})$ என்பது, உள்ளீடு $y_0$,
\dots,~$y_{n-1}$ க்கு முன்னால் $a_0$, \dots,~$a_{m-1}$ ஐ
உள்ளீட்டுச் சரத்தில் இணைத்து,~$e$ ஐ இயக்கும் டியூரிங் பொறி.
அப்போது ஒவ்வொரு $s^m_n$ உம் உரிய டியூரிங் பொறிக்கான குறியீட்டைக்
காணும் முதன்மை மீள்வரையறைச் சார்பு மட்டுமே.}{}
\end{explain}

\end{document}
